\documentclass[11pt,a4paper]{article}
\usepackage{../shared/luxcover}
\usepackage{../shared/proposal}
\usepackage[utf8]{inputenc}
\usepackage[T1]{fontenc}
\usepackage{amsmath}
\usepackage{amssymb}
\usepackage{booktabs}
\usepackage{listings}
\usepackage{xcolor}
\input{../shared/lstlang}

\newcommand{\code}[1]{\texttt{\small #1}}
\newcommand{\prov}[1]{{\color{muted}\footnotesize\sffamily[#1]}}
\setlength{\emergencystretch}{3em}
\sloppy

\title{Morpheus on Lux}
\author{Zach Kelling\\\textit{Lux Industries, Inc.}\\\texttt{research@lux.network}}
\date{2026}

\begin{document}
\luxcoverpage

{\small
\noindent This document is the master record of a read-only evaluation of the Morpheus
inference marketplace, conducted 2026-08-26 through 2026-08-29 against all 62 repositories in
the \code{MorpheusAIs} GitHub organization and against live on-chain state on Ethereum,
Arbitrum, and Base. No repository here is a fork of ours; no change was pushed anywhere. It is
written for the Morpheus core team and their investors. Numbers carry a provenance label:
\prov{verified} read line by line from source or chain, \prov{derived} computed from verified
inputs, \prov{illustrative} assumed for a worked example and not measured, \prov{pending} not
yet true and conditional on a stated event.\par}

\vspace{1.2em}
\tableofcontents
\clearpage

%======================================================================
\section{Executive summary}
%======================================================================

Morpheus describes itself as decentralized, verified, and agentic. The deployed code is
careful engineering at the layers where most projects are weak, and it does none of those
three things at the layer where it counts. This is not an accusation. Every gap named below is
disclosed in Morpheus's own code, comments, or configuration, which by any honest standard
makes it immaturity rather than deception. The point of the evaluation is narrower and more
useful: each of the words in the marketing already maps to a system we run in production on
Lux, Hanzo, or Zoo, so the distance between what Morpheus says and what Morpheus enforces is a
distance we can close with code that already exists.

The core finding is one sentence. A provider is paid \code{pricePerSecond} times elapsed
seconds (\code{SessionRouter.sol:264} \prov{verified}), with no reference in the payment
calculation to the model served, the tokens produced, or whether any inference occurred, and
no slashing anywhere in the contract tree (a \code{grep} for \code{slash|penalt|forfeit} over
the diamond returns nothing \prov{verified}). The receipt the provider signs already carries
\code{inputTokens} and \code{outputTokens} (\code{morrpcmessage/abi.go:12-20} \prov{verified}),
and the contract decodes five of its seven fields and drops both counts
(\code{SessionRouter.sol:236-240} \prov{verified}). Settlement is therefore fakeable, and every
economic mechanism built on top of it, including the buyback that is supposed to give the token
value, inherits that defect.

The economics around it are legible and unflattering. MOR trades at \$1.95, a market
capitalization near \$17M, with 24-hour volume of \$8,905 \prov{verified}. Tradeable DEX
liquidity is \$1,411,470, of which the protocol itself owns 88.8\% \prov{verified}. Real
external revenue, the Lido and Aave yield skimmed from deposits, was about \$476K over the
trailing four quarters and is falling \prov{verified}; MOR emission over the same period is
valued near \$5M to \$6M \prov{verified}. Emission outruns revenue by roughly ten to one. The
schedule that produces it (14,400 MOR/day declining to zero in 2040, 42,000,000 terminal, no
premine \prov{verified}) is the best property the token has, and the buyback that answers the
dilution (about \$511K over the trailing year \prov{verified}) is a sound loop that is simply
outrun.

The thesis follows directly. Keep MOR thin. Morpheus does not need to become a chain to become
real; it needs the settlement it already advertises and a source of value that scales with
usage rather than with deposits. The heavy machinery, post-quantum consensus, verified
inference, threshold custody, a learned router, order-book liquidity, self-repaying credit,
governance with a timelock, already runs on our stack. MOR taps it. It does not rebuild it.

The path is three phases with gates between them. Phase 0 stays on Base with no migration and
is reversible: four one-transaction fixes, a verified-settlement adjunct, and a demand source.
Phase 1 adds governance, an order-book venue, a bank rail, self-repaying credit, and a native
bridge. Phase 2, only if volume justifies it, moves settlement onto a Lux L1 with MOR bridged
as gas, introduces a value token, and cuts emission. Each phase turns on demonstrated on-chain
revenue, not on a promise.

The discipline that governs this document is stated once and kept throughout. We do not claim
what code does not enforce. That rule applies to our own stack, and Section~\ref{sec:heal}
carries our punch-list in the first person, unedited, because a pitch that Morpheus should
adopt our stack ``to become real'' is dishonest until our own stack meets the standard we are
selling. Where we are early, we say early.

%======================================================================
\section{Findings}
\label{sec:findings}
%======================================================================

Four categories. Each finding carries an identifier, the evidence in source or on chain, and a
severity we have deflated rather than inflated. The detailed write-ups live in the security
assessment and the decentralization study; this section is the verified summary and the credit
where it is due.

\subsection{Security and control}

\textbf{CTL-1, Critical.} A 5-of-9 multisig with no timelock has unilateral, instant, unbounded
control over mint, upgrade, and treasury. Every capital contract is UUPS with
\code{\_authorizeUpgrade} gated only by \code{onlyOwner} (\code{DistributionV6.sol:776},
\code{DepositPool.sol:761}, and the same pattern across the tree \prov{verified}). The token
has no supply cap; \code{mint} requires only \code{isMinter[msg.sender]} and \code{updateMinter}
is \code{onlyOwner} (\code{MOROFT.sol:50,61} \prov{verified}). The owner is one Gnosis Safe,
\code{0x1FE04BC15Cf2c5A2d41a0b3a96725596676eBa1E} \prov{verified}. Five keyholders acting
together can mint arbitrary MOR, upgrade the pool that custodies about \$20M of staked stETH
\prov{verified}, or retarget the emission stream, in a single block, with no delay in which a
depositor could observe and withdraw. Every other high-severity item reduces to ``the owner
does X.'' This one is the root, and it is fixable in one transaction: route ownership through a
timelock.

\textbf{CTL-2, High.} The owner directs 52\% of emission with no on-chain contribution check
(\code{manageUsersInPrivatePool}, \code{DistributionV6.sol:175}, \code{onlyOwner}
\prov{verified}); the contributor-to-weight mapping lives in a document, not on chain.

\textbf{SUP-1, High.} The 42,000,000 figure is a curve parameter in the reward math, not a
token invariant. The token itself has no ceiling; the same owner can add a minter and mint
without bound. The ``capped supply'' property is enforced by convention, not by the contract.

\subsection{Settlement and verifiability}

\textbf{SET-1, Medium.} Settlement pays for elapsed time and reads nothing about the work
delivered. \code{modelId} rides on the bid and never enters the bill
(\code{SessionRouter.sol:264} \prov{verified}). A provider is paid the same for the advertised
model, a cheaper substitute, or a degraded quantization.

\textbf{SET-2, Medium.} The two signed token counts are dropped at the contract boundary. The
receipt signs seven fields; \code{\_extractProviderReceipt} decodes five
(\code{SessionRouter.sol:236-240} \prov{verified}). The gap is exactly \code{inputTokens} and
\code{outputTokens}. The protocol discards, at no saving, the one authenticated work record it
already collects. The fix reads two words.

\textbf{SET-3, Medium.} The ranked party signs its own ranking input. \code{tps} and
\code{ttft} are accumulated by the provider's own \code{proxy-router} and feed a five-term
scorer, none of which measures answer correctness (\code{rating/scorer\_default.go:21-25}
\prov{verified}).

\textbf{SET-4, Medium.} Receipt validation checks the signer, not the counterparty, and has no
self-dealing check; nothing requires consumer and provider addresses to differ
(\code{SessionRouter.sol:537} \prov{verified}). Bounded by the 1\% daily budget, so a leak, not
a drain.

\textbf{SET-6, Informational.} Attestation proves an enclave ran, not which weights ran or that
a token reached the user. This is a property boundary, stated so the threat model is explicit:
attestation is not inference verification.

\subsection{Economics}

\textbf{ECO-1, Medium.} Emission outruns real revenue by roughly ten to one, and the gap widens
as price falls. Real yield was about \$476K trailing year and falling steeply, from \$3.24M in
Q1 2024 to \$63K in Q3 2026 \prov{verified}; emission is valued near \$5M to \$6M/year
\prov{verified}. A lower token reduces the deposit incentive, which reduces yield, which reduces
the buyback that supports the token. That feedback runs the wrong way.

\textbf{ECO-2, Medium.} Reselling third-party API quota is a documented provider path
(\code{docs/providers/resale/reselling-venice.mdx} \prov{verified}), and nothing on chain
distinguishes a resold key from GPU compute.

\textbf{ECO-3, Low.} No slashing exists; provider downside is a temporary dispute-hold. A
misbehaving provider risks time, not capital.

\textbf{ECO-4, Informational.} The yield-funded buyback (about \$511K trailing year, 75\% of
captured yield to buy-and-burn, buy-and-lock, and protocol-owned liquidity \prov{verified}) is
well designed and sound. Its only problem is scale.

\subsection{Decentralization and scale}

Decentralization is graded as a vector over six powers (upgrade, custody, halt or censor,
governance, exit, liveness), where the real score is the lowest load-bearing power, not an
average. Morpheus today scores 0 on upgrade, custody, halt, and governance: one Safe can swap
the compute diamond in one block, self-add as minter, recut settlement, and there is no
on-chain vote surface (\code{MOROFT.sol} is a LayerZero OFT with no vote checkpoints; the MRC
process is 86 Markdown files that bind nothing, dormant since 2025-11-17 \prov{verified}). The
honest counterweight is in Section~\ref{sec:heal}: our own stack does not yet clear this bar
either, and we say where.

On scale, the hosted gateway funds every Web2 caller from one shared consumer wallet; after a
nonce conflict on a session open or close the path serializes, and throughput clears a few
sessions per second \prov{verified}. That is an architectural ceiling, not a tuning knob.

The concentration and liquidity facts, read from chain on 2026-08-29 and presented as risks to
depositors rather than as accusations of intent. Existing supply is 9,286,960 MOR (Base 63\%,
Arbitrum 36\%, Ethereum under 1\%) \prov{verified}. Price was \$1.95, market capitalization
\$16.96M, 24-hour volume \$8,905, down 63.9\% on the year and 98.6\% from the all-time high
\prov{verified}. The Compute Reserve multisig holds roughly 2.87M MOR, 30.9\% of all MOR
\prov{verified}; the top ten holders on Base are 93.6\% of Base supply \prov{verified}. Tradeable
DEX liquidity is \$1,411,470, of which the protocol owns 88.8\%, with a protocol-owned-liquidity
Safe holding 99.98\% of the main pool's active liquidity \prov{verified}. Daily volume is 0.63\%
of liquidity; a \$100,000 sell equals about eleven days of total market volume \prov{derived}.
Price discovery and exit both depend on liquidity the protocol controls. This is a fragility any
depositor should be able to see stated plainly, and it is one a governance and liquidity plan
(ADOPT-03, ADOPT-06) should address.

\subsection{What Morpheus built well}

This is not throat-clearing; these are the parts we would keep if we were building on their
code, and Section~\ref{sec:keep} says so in detail.

\begin{itemize}
  \item \textbf{The attestation engineering is serious.} 5,439 lines of Go bind a TDX
    \code{report\_data} to a TLS certificate digest and check an NVIDIA \code{eat\_nonce},
    with RTMR3 computed in CI against a signed manifest \prov{verified}. Most marketplaces
    claim confidential compute and ship nothing; Morpheus shipped a working measurement path.
  \item \textbf{The treasury bifurcation did its job.} Consumer-funded and treasury-funded
    payments take different code paths, and the 1\% daily budget cap bounded the July 2026
    stipend leak to roughly 25,000 MOR/day rather than draining the reserve \prov{verified}.
  \item \textbf{The receipt already carries the work record.} Seven signed fields, including
    both token counts \prov{verified}. The harder half of verifiable settlement, getting a
    signed structured record off the provider, is already built.
  \item \textbf{Accounting is receipt-derived and the audit trail is real.} No premine; the
    \code{Docs/Security Audit Reports/} directory holds genuine Code4rena, Zenith, and
    Renascence reports \prov{verified}. This is a project that paid for adversarial review and
    kept the receipts.
\end{itemize}

The Ponzi accusation, tested directly, is refuted. Depositor principal is withdrawable 1:1
(\code{DepositPool.sol:265} \prov{verified}) and never redistributed; rewards are freshly
minted MOR, structurally a block reward. The defensible statement is ECO-1: an inflationary
reward under-earned by real demand, which is a sustainability question, not fraud.

%======================================================================
\section{The proposal, as a menu}
\label{sec:menu}
%======================================================================

The response is thirteen one-page proposals, the ADOPT series, each standing for a longer
paper and each independently readable. They are organized into four tiers by dependency depth,
not by importance. A tier is a floor: everything in it stands on what is below it. Morpheus can
take one tier, or one paper, or all four. Nothing here is a bundle that fails unless bought
whole.

\begin{center}\small
\begin{tabular}{@{}llll@{}}
\toprule
\textbf{Docket} & \textbf{What it fixes} & \textbf{Tier} & \textbf{Depends on} \\
\midrule
ADOPT-09 & CTL-1 timelock, SET-2/5 decode and bind, SET-4 self-deal & 1 & none \\
ADOPT-02 & SET-1/2/3, ECO-3: pay for work, not time & 1 & 09 \\
ADOPT-07 & ECO-1 demand, ECO-2 resale, SET-6 attested supply & 1 & 02 \\
\midrule
ADOPT-03 & CTL-1/CTL-2: timelock, vote surface, owner emission & 2 & 09 \\
ADOPT-08 & ECO-1/ECO-4: buyback fed by volume, not deposits & 2 & 02 \\
ADOPT-05 & fiat on-ramp: accounts and cards & 2 & none \\
ADOPT-06 & thin liquidity: order book and price discovery & 2 & 03 \\
ADOPT-04 & idle \$20M stETH, exit pressure: LMOR credit & 2 & 06 \\
\midrule
ADOPT-13 & hand-replicated cross-domain schedule: native bridge & 3 & 04 \\
ADOPT-10 & demand loop: reciprocal routing across the group & 3 & 07, 08 \\
\midrule
ADOPT-01 & no chain, validator set, or consensus of its own & 4 & 03, 13 \\
ADOPT-12 & keep MOR thin: tap, do not rebuild & 4 & 01 \\
ADOPT-11 & few sessions per second: settlement throughput & 4 & 01 \\
\bottomrule
\end{tabular}
\end{center}

\textbf{Tier 1, on Base, no migration, reversible.} The four one-transaction fixes protect
depositors and cost nothing structural (ADOPT-09). Verified settlement makes the meter read the
token counts the receipt already signs (ADOPT-02). A demand source and attested supply make the
meter worth reading (ADOPT-07). This tier closes CTL-1, SET-1 through SET-5, ECO-3, and half of
SET-3 without moving a single asset off Base.

\textbf{Tier 2, value and governance.} A timelock and a vote surface move control off the bare
multisig (ADOPT-03). A fee on verified volume funds the buyback from usage instead of deposits
(ADOPT-08). An order-book venue deepens a book that is 88.8\% protocol-owned (ADOPT-06).
Self-repaying credit puts the idle stETH and MOR to work without a forced sale (ADOPT-04). A
bank rail brings fiat-native consumers in (ADOPT-05, the smallest dependency in the set).

\textbf{Tier 3, cross-domain.} A native MPC bridge replaces the schedule that is replicated by
hand across the capital and compute domains (ADOPT-13). A reciprocal routing loop turns the
group's own inference demand into MOR volume (ADOPT-10).

\textbf{Tier 4, a chain of its own, only if warranted.} A Lux L1 or L2 gives Morpheus the
consensus, validator set, and finality it does not have today, inherited rather than built
(ADOPT-01). MOR stays thin on top of it (ADOPT-12). Settlement throughput stops being a few
sessions per second (ADOPT-11). This tier is optional and gated on volume; Section~\ref{sec:path}
states when it is worth the migration and when it is not.

%======================================================================
\section{Where MOR goes long term}
\label{sec:mor}
%======================================================================

MOR today is one token doing three jobs. It is emitted on a schedule, escrowed to open
sessions, and paid to providers. Any token-evolution story has to fix settlement before the
emission problem is worth discussing, because every value mechanism sits downstream of the
meter.

\subsection{Two tokens, and what gives the second one value}

Split the roles. MOR stays the inflationary budget-and-rewards unit: emitted, spent as gas and
escrow, paid to providers. xMOR becomes the low-velocity store of value. The question that
decides whether xMOR is real or decorative is simple: what claim does an xMOR holder have that
a MOR holder does not?

\textit{Burn-and-mint} is the weaker construction. Burn MOR to mint xMOR at some ratio. If the
ratio is fixed and reversible, xMOR is a wrapper and accrues nothing; if one-way, it is MOR
minus an exit, worth less rather than more, unless xMOR separately carries a right. The mint
ratio is a governance parameter, and a governance-set conversion between two pools of value is
exactly the surface the same multisig that directs 76\% of emission can capture. Burn-and-mint
values throughput, not savings.

\textit{Share-token}, in the xSUSHI style, is sounder. Stake MOR into a pool; xMOR is the
share. Protocol fees buy MOR and add it to the pool, so each xMOR redeems for a growing
quantity of MOR. Value is an endogenous cash-flow claim, not a parameter: the xMOR/MOR ratio
rises monotonically with fees collected, and there is no conversion ratio for anyone to set.
Its one trust point is that fees actually reach the pool, which every honest design shares. Its
only leakage is timing, deposit before a distribution and withdraw after, and streaming the
reward over a vesting window collapses that arbitrage.

Pick the share-token, and state the thing the client should hear. If fees buy MOR and burn it,
every MOR holder benefits equally and no second token is needed at all. The only defensible
reason to mint xMOR is velocity separation, keeping the budget churn from diluting the balances
people hold for value. That is a real reason. It is not a magic one.

\subsection{Burning does not create value}

Say it plainly. Price is market capitalization divided by supply, so shrinking supply lifts
price only if market capitalization holds. Burning adds nothing to market capitalization. A
deflationary token with no demand is the same pie cut into fewer slices. Burn is real
value-accrual in exactly one case: when it is financed by fees that real users paid for real
usage. Then the fee is the value entering the system and the burn is merely how that value
reaches holders, as scarcity instead of as a dividend. A burn financed by emission or token
sales is circular: emit, sell, buy back, burn, net zero minus slippage.

This is where our own discipline bites. Settlement is fakeable today, so a fee on unverified
volume is a fee on a number a provider can fabricate, and a self-dealing provider will wash
inference to drive a burn that lifts a token he holds. The burn is only as honest as the meter
behind it. Verified settlement (ADOPT-02, and \code{lux/mord} at about 3,400 verified
sessions/sec/core \prov{verified}, a settlement-rate figure, not network TPS) is the
precondition, not a companion feature.

\subsection{Utility is the engine}

Volume comes from MOR being the unit of account for agentic compute: agent budgets, MCP-server
calls, inference and compute generally. Every call spends MOR, and a fee on each call funds the
buy-and-burn. Define the variables. Let $E$ be emission in MOR/day, $A$ the number of
fee-bearing calls/day, $c$ the average revenue per call in USD, $f$ the fee fraction routed to
buy-and-burn, and $P$ the MOR price in USD. MOR burned per day is $B = fAc/P$. Net supply change
is $E - B$, and the system is net-deflationary when
\[
  B > E, \qquad\text{that is,}\qquad Ac > \frac{EP}{f}.
\]
The right side is the emission's daily dollar value divided by the fee rate. Everything is
legible from there.

A worked example, verified inputs and an illustrative conclusion. Emission's daily dollar value
is $EP \approx 14{,}400 \times \$1.95 \approx \$28{,}080$/day \prov{derived}. At $f = 1\%$,
net-deflationary needs gross fee-bearing volume $Ac > \$2.81$M/day, about \$1.0B/year
\prov{illustrative}. At $f = 5\%$, above \$205M/year; at $f = 10\%$, above \$102M/year. Against
today's roughly \$476K/year of revenue, which is Lido and Aave yield rather than usage, that is
the true size of the ask. The assumptions (constant price, fee taken in USD-equivalent and
fully burned, all volume verified) do not hold exactly; the equation shows the shape, not a
forecast.

\subsection{Cutting emission moves the same inequality}

$E$ is not a protocol constant needing a hard fork. It is set by owner-controlled pool
configuration today (\code{manageUsersInPrivatePool}, \code{DistributionV6}, \code{onlyOwner}
\prov{verified}). Lower $E$ and the crossover bar $Ac > EP/f$ falls proportionally. Halve
emission to 7,200/day and the 1\% requirement halves to about \$510M/year; cut it tenfold and it
falls to about \$102M/year \prov{illustrative}. Emission reduction and a usage-driven burn are
the same lever pushed from two ends. The honest footnote is that the same owner power that can
cut emission is the centralization CTL-1 and CTL-2 flag; the ability is real, and so is the
risk, which is why ADOPT-03's timelock precedes any emission change.

\subsection{The ICO funds the transition, not the demand}

On \code{lux.exchange}, conditional on SEC approval of the proposed regulation \prov{pending}.
Its job is price discovery, treasury, and protocol-owned liquidity, not a substitute for
demand. The float is thin (depth about \$1.4M, 88.8\% protocol-owned, 24h volume about \$8.9k
\prov{verified}), and a raise deepens the book and funds the build-out of the fee-bearing
utility. It does not manufacture usage. The raise should clear only into demonstrated on-chain
revenue. Selling the token before the meter is honest and the utility exists is selling the
thing we just called theater.

\subsection{What can go wrong}

Two-token models add complexity and a conversion surface someone can capture; prefer
buy-and-burn on MOR unless velocity separation earns the second token. A burn with no demand is
theater, and a burn on unverified volume is worse. Emission cuts take income from the recipients
of the 76\% owner-directed emission, who will resist. The ICO needs real demand first or it
prices thin float into a spike that round-trips. And every appreciation claim here is downstream
of one unshipped thing: verified settlement. Ship the meter, then the fee, then the burn, then
the token that holds its value. In that order.

%======================================================================
%======================================================================
\section{Listing and liquidity}
\label{sec:listing}
%======================================================================

ADOPT-06 named the venue problem and the order-book fix: a book that is 88.8\% protocol-owned,
where a \$100k sell is eleven days of volume, and where a two-sided book replaces a passive
reserve \prov{verified}. This section is the layer above the venue. It is about how MOR reaches a
liquid market, what kind of market that is, and how far that market extends. Three moves, ordered
by how much review each needs: list MOR, bridge and lend against it, and carry the broader
cross-listed market above it.

\subsection{MOR as a first-class Lux asset}

MOR is \code{Offering: unregulated} in the asset taxonomy \prov{verified: asset-class.ts}, so
none of this touches the securities surface. Three wirings, all crypto, all shippable now:

\begin{itemize}
\item \textbf{Featured listing.} The interface's token list pins the ecosystem set --- MOR
      alongside AI, ZOO, and LUX --- at the top of the surface a user sees first. This is a
      one-file change at a single chokepoint, and the four are defined once as a typed
      \code{FeaturedToken} list whose class field refuses a security by construction.
      \prov{verified: pkgs/lx TokenSelector suggested-tokens hook; \code{FeaturedToken} in pkgs/exchange/src/featured.ts}
\item \textbf{Native bridge.} MOR is a LayerZero OFT on Ethereum, Arbitrum, and Base
      \prov{verified}. Lux's own bridge already carries Arbitrum and Base, so MOR enters the
      set it already spans. Adding MOR is asset registration --- three token entries plus an
      on-chain allowlist --- and an L-token wrapper; no new chain integration.
      \prov{verified: lux/bridge token registry, chain manifests, and \code{allowedTokens}}
\item \textbf{LMOR.} MOR as collateral in \code{lux/liquid} self-repaying lending, so a holder
      borrows against MOR without a forced sale --- the ADOPT-04 move applied to MOR itself, and
      the concrete form of ``put the treasury's idle assets to work.'' Two paths: a Morpho
      market is one transaction and gives overcollateralised MOR credit today; the self-repaying
      market is the ADOPT-04 design and needs a yield-bearing MOR to retire the debt.
      \prov{verified: lux/liquid --- Morpho \code{createMarket} vs the Alchemix-style Liquid market}
\end{itemize}

\noindent What this needs from Morpheus: the MOR OFT address on each chain, and confirmation of
the AI, ZOO, and LUX addresses. Nothing here moves a Morpheus contract or needs the 5-of-9.

\subsection{The interface is a renderer, not a broker}

The venue above carries more than MOR, and the model that lets it do so cleanly is already the
one our interface is built on. It converts a user's chosen parameters into a transaction the user
signs in their own wallet. It never holds a key, never executes or settles, never routes an
order, and is paid only by the user \prov{verified: COMPLIANCE.md}. Under the SEC Division of
Trading and Markets staff statement of 13 April 2026 (File No. 4-894), a provider of such an
interface --- a Covered User Interface Provider --- is not a broker or dealer under Exchange Act
section~15, provided it meets the statement's enumerated conditions. Our \code{COMPLIANCE.md} is
written against that statement directly \prov{verified}.

Stated plainly, because a design built on it inherits its expiry: the statement is staff views
with no legal force, withdrawn 13 April 2031 absent Commission action, and Regulation Crypto
Assets (Release 33-11434) is a proposal, not law \prov{verified}. What follows is designed
against both; it is not compliance with either.

\subsection{Three things such an interface lists}

The live Uniswap interface is the reference. It carries, in one surface:

\begin{enumerate}
\item \textbf{Ecosystem tokens, featured.} For LX, the set above: MOR, AI, ZOO, LUX.
\item \textbf{A launchpad long tail, aggregated.} Uniswap's Launches surface --- an ``all
      launchpads'' filter, a trending list, ``Pools is live: trade new Robinhood Chain tokens''
      --- indexes tokens as launchpads mint them. The first target for LX is pools.trade on
      Robinhood Chain. Aggregation routes to the token's liquidity on its \emph{source} chain; it
      does not bridge the token onto Lux and does not settle a security. In \code{lux/dex} the
      gateway already registers providers by priority; a launchpad is one more route-only
      provider, and adding the next one is a new registry entry, nothing else.
      \prov{verified: lux/dex gateway provider registry; a route-only \code{LaunchpadProvider}}
\item \textbf{Tokenized equities, shown globally and geo-gated.} Uniswap lists NVDA under
      Stocks, ``Issued by Robinhood,'' beside Coinbase, Xstock, Bstocks, and Ondo wrappers of the
      same underlying, and blocks the swap where the issuer's wrapper requires it: ``NVDA
      unavailable in your region'' \prov{verified: live interface}.
\end{enumerate}

\subsection{What we have, and what is switched off}

Two codebases carry this lineage. \code{lux/exchange} is a debranded fork of the Uniswap
interface; \code{lux/exchange2} is a ground-up, securities-native rewrite \prov{verified: git}.
Neither yet renders the Launches or Stocks surfaces; the histories are disjoint, so enabling them
is a feature-directory port, not an upstream merge \prov{verified}. The tokenized-equity taxonomy
already exists --- \code{AssetClass.STOCK} and an \code{Offering} union spanning \code{reg\_cf},
\code{reg\_a\_plus}, \code{reg\_s}, \code{reg\_d}, and \code{rule\_144a}
\prov{verified: asset-class.ts} --- but the render is off: \code{Features.publicSecurities}
defaults false, \code{canTrade} has no caller, and no equity appears in any token list. The
taxonomy exists; the render does not \prov{verified: equities.md, features.ts, regulated-provider.ts}.

The binding constraint is not interface code. Uniswap's Launches and Stocks read its hosted
data-api, which does not serve Lux chains. The real work is the data plane: index launchpad and
tokenized-asset state for Lux, or aggregate it through the \code{lux/dex} gateway, which is
that data plane. \prov{verified: the lux/dex gateway connector is the source}

\subsection{Two boundaries that do not move}

\begin{description}
\item[Interface versus issuer.] Listing someone else's tokenized security is the
      Covered-User-Interface posture above. Issuing our own offering --- a MOR raise, or any Reg
      CF, Reg A+, Reg S, or Reg D placement --- is the issuer or portal role, a different and
      heavier one. Both are modeled in the \code{Offering} union; the ICO in
      Section~\ref{sec:mor} is the issuer path, gated on the proposed regulation
      \prov{verified}.
\item[Aggregation versus bridging.] Featuring MOR, an OFT we bridge natively, is not the same as
      carrying a tokenized equity. An equity token keeps its issuer's non-US wrapper wherever it
      trades; the interface mirrors that restriction, it does not shed it by rendering. So the
      equity surface is shown globally, gated in the United States pending the rules now being
      finalized, and it ships only behind counsel. The flag flip is the client's, not ours ---
      the render is off by a flag with no reader today, and the region block is net-new: only a
      UK information banner exists, so a fail-closed US block is a control to build, not toggle.
      \prov{verified: features.ts, regulated-provider.ts, the UK-only region banner}
\end{description}

\noindent MOR is the shippable core today. The launchpad tail follows as the gateway connector
lands. The equity surface is a built-but-gated capability whose switch is a legal decision, not
an engineering one. One interface, three depths of review.

\section{What we bring}
\label{sec:bring}
%======================================================================

This is the catalogue of what is available to adopt. Every entry is a real repository or paper
on disk. Maturity is one of \textbf{verified live} (running in production on Lux, Hanzo, or Zoo,
not yet on Morpheus), \textbf{built, not deployed} (code with tests, not activated on a live
chain), or \textbf{research} (paper or proof, not shipped). One caveat is kept throughout:
``gives Morpheus'' means available to adopt, not already wired in. Nothing here settles a
Morpheus session today.

\subsection{Consensus and post-quantum}

\textbf{Quasar 3.0} (\code{lux/consensus}, \code{lux/quasar}; LP-020). DAG-native BFT that
finalizes a block with one \code{QuasarCert} aggregating three lanes: a BLS12-381 classical fast
path, a Corona Module-LWE threshold signature, and a Z-Chain verifier rollup. \textbf{Verified
live} (LP-020 states activation on the Lux primary network 2025-12-25). For Morpheus, this makes
the thin-MOR thesis concrete: a Morpheus L1 for session and settlement state can rent classical-
fast and PQ-durable finality on the same certificate rather than build consensus.

\textbf{Pulsar / Corona / Magnetar threshold family} (\code{lux/pulsar}, \code{lux/threshold};
LP-073). Lattice threshold signatures with a lifecycle framework for permissionless validator
sets; Pulsar (threshold ML-DSA) is an active NIST MPTC submission. \textbf{Built, not deployed}
as a reusable primitive. For Morpheus, a threshold-sealed close-of-session receipt signed by a
provider set rather than one hot secp256k1 key, which addresses the single-hot-key custody
problem directly.

\textbf{Unified PQ transport} (\code{lux/papers} LP-202). One UDP port carrying QUIC over ZAP,
an X25519MLKEM768 hybrid KEM in the TLS 1.3 handshake, hybrid classical and ML-DSA-65 peer
identity. \textbf{Research} (specification). For Morpheus, the transport fix for the plaintext
prompt path: prompts currently cross \code{:3333} in cleartext, and this is the drop-in PQ
session layer.

\textbf{P3Q unified-slot verifier} (\code{lux/precompile/p3q}; LP-219). One EVM precompile slot,
one ABI, a leading \code{kind} byte dispatching per-family threshold-signature verifier cores.
\textbf{Built, not deployed}. For Morpheus, any contract, a settlement diamond included, can
verify PQ threshold certificates on-chain without a bespoke precompile per scheme.

\subsection{Verified inference and A-Chain}

\textbf{Inference precompile 0x0303} (\code{lux/precompile/inference}). A deterministic int8
transformer as an EVM precompile, pure Go with CGO=0, so every validator recomputes it
byte-identically; GPU kernels are a KAT-proven byte-equal accelerator. \textbf{Built, not
deployed}. For Morpheus, the missing settlement input: a deterministic small-model verification
path is one honest way to bind payment to work for models small enough to reprice on-chain.

\textbf{Attestation precompile 0x0301} (\code{lux/precompile/attestation}). Stateless TEE-quote
verification (NVTrust GPU, TPM/SGX/SEV-SNP/TDX), chained to an embedded vendor root at the
quote's own timestamp, fail-closed, no mutable device registry. \textbf{Built, not deployed}.
This is the primitive behind Morpheus's own live TEE work (RTMR3, SecretVM), but as a
consensus-checkable precompile rather than an off-chain check.

\textbf{A-Chain inference bridge} (\code{lux/precompile/aivmbridge},
\code{lux/precompile/modelregistry}). Two operations, submit-inference-intent
(\code{modelSpecHash}, \code{promptHash}) and verify, plus a model registry precompile.
\textbf{Built, not deployed}. For Morpheus, a \code{modelId} that is a registered, hashed model
spec rather than a line in a local JSON file, which makes the resale substitution ECO-2
documents detectable.

\textbf{Lux AI / A-Chain} (\code{lux/ai}). OpenAI-compatible decentralized inference network,
Q-Chain shared attestations, cross-network minting. \textbf{Built, not deployed} (our
punch-list fixed a \code{lux/ai} unauthenticated RPC on 2026-08-29; see
Section~\ref{sec:heal}). For Morpheus, an existing inference-settlement chain to tap under
thin-MOR rather than growing Morpheus into one.

\textbf{Zoo proof-of-inference} (\code{zoo/proofs}). A Freivalds matrix-product verifier over
$\mathbb{F}_p$ ($p = 2^{61}-1$) that catches a fabricated product with probability at least
$1 - 1/p$ per challenge, zero false rejection over an exact integer accumulator, wired to an
activation trace so an operator cannot mint against an output it never computed.
\textbf{Research} (proofs, some Lean-formalized). For Morpheus, a cheap probabilistic check that
a claimed forward pass actually happened.

\textbf{mord} (\code{lux/mord}). A Lux L1 that runs an inference marketplace as chain state,
separating capacity, work, and identity into three settlement quantities instead of Morpheus's
one \code{pricePerSecond * elapsed}. \textbf{Built, not deployed} (our POC: 40 tests plus
adversary tests, benchmarked at about 3,400 verified sessions/sec/core; not deployed, not
audited; that number is settlement-rate, not network TPS). The reference design for what
un-fakeable settlement looks like on their own market shape.

\subsection{Agentic runtime}

\textbf{Enso learned router} (\code{hanzo/enso}, ledger at
\code{hanzo/ai/object/routing\_event.go}). Classifies each request, prices it, routes to the
best model across enabled providers, learns from in-app feedback and an LLM-as-judge signal,
routes in microseconds on CPU. The ledger stores task label, routed model, token counts,
latency, cost, an optional opaque feature vector, and a scalar reward, and stores no prompt text
or hash. \textbf{Verified live} (powers \code{model=auto} on Hanzo Cloud). For Morpheus, a
router that improves on a time-priced market and is the deployed proof that confidential routing
is buildable: it consumes a low-dimensional statistic plus a scalar outcome, not content.

\textbf{Hanzo agent SDK, agents, runtime} (\code{hanzo/agent}, \code{hanzo/agents},
\code{hanzo/runtime}). Multi-agent orchestration over an MCP tool surface, role and workflow
agents, and a sandboxed execution runtime for AI-generated code. \textbf{Built and verified
live} (runtime ships releases). For Morpheus, the agentic execution layer that sits on sessions
as the substrate for agent turns.

\textbf{MCP tool surface} (\code{hanzo/chat}, \code{mcp\_\_hanzo\_\_*}). Model Context Protocol
tool use wired through \code{api.hanzo.ai/v1}. \textbf{Verified live}. Standard tool-calling for
provider-side agents without inventing a protocol.

\textbf{Hanzo Dev} (\code{hanzo/dev}). A terminal coding agent, sandboxed repo edits, models
over the metered gateway. \textbf{Verified live}. A reference consumer of metered per-session
inference.

\subsection{Confidential compute}

\textbf{FHE precompile} (\code{lux/precompile/fhe}). CKKS via the pure-Go \code{luxfi/lattice},
threshold decryption through a T-Chain (LP-333). \textbf{Built, not deployed}. For Morpheus,
compute over consumer data where the operator should not read inputs.

\textbf{ZK, STARK-FRI, and anchor precompiles} (\code{lux/precompile/zk},
\code{lux/precompile/starkfri}, \code{lux/precompile/anchor}; LP-221). Proof-verification and
commitment precompiles, plus an anchor precompile that stores a Merkle root of CRDT state with
per-app height monotonicity. \textbf{Built, not deployed}. The anchor primitive is exactly what
the two-economy problem needs: Morpheus runs capital on Ethereum and compute on Base with a
schedule replicated by hand (\code{SessionRouter.sol:58} says so \prov{verified}), and anchoring
lets one domain verify the other's state root instead of trusting a manual copy.

\textbf{MPC custody stack} (\code{lux/threshold}, \code{lux/kms}, \code{lux/mpc}). Universal
threshold signing across many chains (CGGMP21 ECDSA, FROST Schnorr), MPC-backed KMS with per-org
namespaces. \textbf{Built and verified live} (the canonical Lux KMS). For Morpheus, splitting the
one provider hot key into capacity, work, and identity authorities with different blast radii.

\subsection{Settlement, bridge, governance}

\textbf{Lux primary settlement and MOR OFT} (\code{lux/node}, LayerZero OFT). MOR is already a
live LayerZero OFT on Ethereum, Arbitrum, and Base. \textbf{Verified live}. For Morpheus,
settle on a chain with real finality instead of application-layer storage on one Base diamond.
The honest counterweight (LP-314) is that a sovereign L1 buys nothing worth the migration for
much of the market, so this is an option, not a mandate.

\textbf{Lux Bridge} (\code{lux/bridge}). MPC cross-chain bridge, embeddable SDK,
\code{api.bridge.lux.network}. \textbf{Verified live}. Native movement of MOR and stake across
the capital and compute domains without a hand-replicated schedule.

\textbf{l2-template and Quasar-rooted L1} (\code{lux} L2 template, \code{lux/evm}). An L2 stood
up in minutes, rooted to Quasar PQ finality. \textbf{Built}. The minimal-effort route to a
Morpheus-owned chain that inherits PQ finality rather than renting rollup execution.

\textbf{Lux DAO and vote} (\code{lux/dao}, \code{zoo/vote}). Governance contracts plus a UI,
white-labeled per org. \textbf{Verified live} (\code{zoo/vote} deployed). An on-chain governance
surface for parameter changes and treasury actions. Stated plainly: adopting a DAO UI does not
by itself add a timelock or bound the mint; those are contract changes.

\subsection{Fintech}

\textbf{LX order-book DEX} (\code{lux/dex}). Pure-Go matching engine, order book, oracle
aggregator, JSON-RPC/WebSocket/gRPC SDKs. \textbf{Verified live} (reference implementation runs
standalone, ships releases). For Morpheus, real order-book depth for MOR against the thin AMM
the protocol itself backstops.

\textbf{Lux Exchange} (\code{lux/exchange}). A Uniswap-forked front end plus a securities and
ICO surface. \textbf{Built}; the securities/ICO path is \textbf{pending} SEC approval of a
proposed regulation, a stated dependency and not a granted approval. Do not represent it as
available.

\textbf{Liquid self-repaying lending} (\code{lux/liquid}; LP-310). Deposit yield-bearing
collateral, borrow a 1:1-pegged synthetic, let yield retire the debt on a time-decay schedule;
curated strategy adapters (Aave, Lido, Pendle, Morpho, Yearn). \textbf{Built, not deployed}. For
Morpheus, an LMOR construction: self-repaying credit against MOR or against the roughly \$20M
stETH \prov{verified} the treasury already holds, so holders borrow without selling into the
\$8.9k/24h \prov{verified} volume.

\textbf{Lux Bank} (\code{lux/bank}, \code{lux/captable}). Accounts, cards, cap-table.
\textbf{Built, not deployed}. A fiat and accounts on-ramp for the hosted-gateway audience,
account-native consumers behind delegated sessions rather than one shared hot wallet.

\subsection{Data and observability}

\textbf{GPU-native EVM and verification} (\code{luxfi/evm}; LP-203). Unified C-Chain and L1 EVM
plus GPU-native transaction verification. \textbf{Built} (EVM ships releases) and
\textbf{research} (the billion-TPS figure is a paper claim, not a measured production number,
labeled illustrative). Batched signature and proof verification throughput for a settlement
chain. As with mord, the benchmark rate is not sustained network TPS.

\textbf{Hanzo gateway, IAM, telemetry} (\code{hanzo/gateway}, \code{hanzo/iam}). The trust
boundary at \code{api.hanzo.ai} (JWT verified against IAM JWKS, client identity rewritten, rate
limits, circuit breakers) and OIDC identity at \code{hanzo.id} (PKCE, SCIM, WebAuthn/MFA).
\textbf{Verified live}. Morpheus's \code{api.mor.org} multiplexes many callers onto a few
on-chain sessions behind spend limits that ship inert and are armed per environment; a
verified-identity gateway with armed limits is the deployed answer.

\textbf{Zoo Gym decentralized training} (\code{zoo/gym}). An OSS LLM fine-tuning framework plus
a protocol paper for ring-topology training across heterogeneous GPUs. \textbf{Built}
(framework) and \textbf{research} (decentralized-training protocol). A provider-side path to
sell training and fine-tuning yield, not only inference, on the same GPU fleet.

\subsection{What this catalogue does not claim}

Kept explicit. Nothing here settles a Morpheus session now; every ``gives Morpheus'' is an
adoption, not a wiring. ``Verified live'' means live on Lux, Hanzo, or Zoo; the only
Morpheus-side live PQ or TEE work is Morpheus's own, credited above. The billion-TPS and
3,400-sessions/sec figures are benchmark and rate numbers, not network TPS. The
\code{lux/exchange} securities and ICO path is pending a regulatory approval that has not been
granted. LP-312 through LP-326 are research papers grounded on Morpheus's deployed contracts;
they are the argument for adoption, not evidence of it.

%======================================================================
\section{Building on their work}
\label{sec:keep}
%======================================================================

A migration that discards the good parts is not an improvement; it is a rewrite with a survivor
bias toward whatever we happened to build. So the honest split is stated first, and it is short.

\subsection{Keep}

\begin{description}
  \item[The attestation stack.] The 5,439-line TDX-to-TLS binding with RTMR3 in CI is real
    confidential-compute engineering. Keep it and give it a consensus-checkable home
    (attestation precompile 0x0301) so what today is an off-chain check becomes a settlement
    input. The local-attestation path in \code{lux/ai/pkg/attestation} removes the external
    NRAS and SecretVM availability coupling (SET-6); adopt it or keep their own, but keep the
    measurement work.
  \item[The bounty and audit design.] Real Code4rena, Zenith, and Renascence reports and a
    well-structured bounty program. Keep the process; it is the standard this evaluation holds
    itself to.
  \item[The treasury bifurcation.] Separate consumer and treasury payment paths with a 1\%
    daily budget cap. It contained the July 2026 stipend leak. Keep the split; the fix that
    day-locked stipend on \code{min(closedAt, endsAt)} is exactly the right shape.
  \item[Receipt-derived accounting.] Stipend size and emission are pure functions, legible and
    independently checkable. Keep the receipt (it already signs both token counts) and read the
    two words the contract currently drops (SET-2).
  \item[The emission schedule.] 14,400/day to zero in 2040, 42,000,000 terminal, no premine.
    The best property the token has. Keep it, and add the invariant cap the token lacks (SUP-1)
    so the property is enforced rather than conventional.
  \item[The community and the token.] MOR stays MOR. The migration keeps the ticker, the
    holders, the OFT deployment, and the schedule. Nothing here asks the community to move to a
    new token.
\end{description}

\subsection{Replace}

\begin{description}
  \item[The settlement calculation.] \code{\_getProviderRewards} pays elapsed time. Replace the
    payable event with a quorum-agreed receipt that commits to the output token count, so the
    fee base is reproducible (ADOPT-02). This is the one seam; everything above and below it
    stays.
  \item[The bare multisig.] A 5-of-9 with no timelock is the root finding. Replace it with a
    Safe plus Governor plus timelock, with a nonzero floor on the delay (ADOPT-03).
  \item[The buyback's fuel.] Deposit yield inverts when price falls. Replace it with a fee on
    verified volume (ADOPT-08), so the loop scales with usage.
  \item[The hand-replicated cross-domain schedule.] Replace the manual copy between Ethereum and
    Base with a native bridge and state-root anchoring (ADOPT-13, anchor precompile).
  \item[The self-signed ranking.] Replace the provider-signed \code{tps}/\code{ttft} scorer with
    a learned router keyed on an independent quality signal (Enso), which is orthogonal to and
    composes with quorum verification (SET-3).
\end{description}

%======================================================================
\section{Heal thyself}
\label{sec:heal}
%======================================================================

This section is first-person and unedited, because the pitch that Morpheus should adopt our
stack to become real is dishonest until our own stack meets the standard we are selling. The
rule it enforces: do not make a claim the code does not yet enforce. Every item below is a
failure in our own tree (\code{luxfi}, \code{hanzoai}, \code{zooai}), verified against source on
2026-08-29. The decentralization study grades six powers as a vector, and ``theater'' is the gap
between a claimed state and the measured one. Each item moves a power off zero or closes a
claimed-versus-measured gap.

\subsection{Update: the on-chain governance path, built and proven (2026-08-30)}

Since the audit above, the governance items (B3, B4, B5) moved from open to \emph{built and
demonstrated}. We stood the full stack up on a local network and exercised the decentralization
end to end, verified on-chain, not described.

\begin{itemize}
\item A one-shot deploy brings up the entire stack --- DEX, emissions, treasury, bridge,
  work-market, and the governance instrument --- with every authority resting on a single 1/1
  Safe: the honest bootstrap state \prov{verified: local net}.
\item The governance instrument is the audited \code{DeployGovernance} wiring, and it comes up
  clean: the deployer renounces the Timelock admin (\code{deployer=false}, \code{self=true}), the
  Governor is the sole proposer, and the Safe is canceller/veto only. No residual deployer or EOA
  power \prov{verified: on-chain reads}.
\item \code{lux.vote} is wired to the live Governor, and an admin dashboard (\code{admin.lux.cloud})
  renders every governable surface with its current owner and a button that hands it from the Safe
  to the Timelock. We fired one such handoff through the 1/1 Safe on-chain --- an AMM factory's
  authority moved Safe~$\rightarrow$~Timelock, transaction status success \prov{verified}.
\end{itemize}

So the mechanism the pitch depends on --- bootstrap under one key, then relinquish every parameter
to token-holder governance --- is no longer a proposal; it runs, and each surface it touches ends
up controlled by the Governor through the Timelock, not by a key. That answers B3, B4, and B5 on
the mechanism: the honest 1/1 bootstrap replaces the 3-of-5-from-one-mnemonic costume, the
deployer-admin timelock is renounced by the correct script, and \code{lux.vote} over a live
Governor replaces a branded front end over null contracts.

\textbf{The claim we make, and the one we withhold.} We claim: Lux has built and demonstrated,
on-chain and end to end, an open on-chain governance system --- \code{lux.vote} --- and the path
from a single bootstrap key to full token-holder control, with the Governor and Timelock live. We
do not claim Lux is \emph{already} fully decentralized on mainnet: today the mainnet parameters
still rest on the bootstrap key, and the network reboot plus the parameter-by-parameter handoff to
the Timelock are the executing step, not a finished one \prov{pending}. Asserting the finished
state before that handoff would be the exact theater we flag in Morpheus. So we state the measured
position: the infrastructure is real and proven, the mainnet handoff is one execution away, and
the same dashboard runs it.

\subsection{The same twelve traps, in our own five systems}

Before the punch-list, the mirror. The five economic and settlement failures we flag in
Morpheus are not uniquely Morpheus's; we ran the same twelve-trap rubric over our own systems.
Status is \code{present} (the defect exists and is reachable), \code{at-risk} (exists on a path
not yet exploitable or a live design gap), \code{avoided} (closed in the reachable path, with
the limit stated), or \code{n/a}. ``papers'' is the LP inference program read as a
specification, so its \code{avoided} means closed by design, not deployed.

\begin{center}\footnotesize
\begin{tabular}{@{}lccccc@{}}
\toprule
\textbf{Trap} & \textbf{lux-ai} & \textbf{mord} & \textbf{hanzo-ai} & \textbf{gov/emission} & \textbf{papers} \\
\midrule
T1 unverified pay for claimed work & at-risk & \textbf{present} & at-risk & \textbf{present} & avoided \\
T2 no slashing of a bond & avoided & \textbf{present} & n/a & \textbf{present} & avoided \\
T3 uncapped mint & at-risk & avoided & n/a & at-risk & present \\
T4 owner upgrade, no delay & n/a & avoided & n/a & at-risk & at-risk \\
T5 no vote surface & n/a & avoided & n/a & at-risk & present \\
T6 owner-directed emission & n/a & n/a & n/a & at-risk & present \\
T7 prompt transport confidentiality & at-risk & n/a & at-risk & n/a & avoided \\
T8 resale indistinguishable & \textbf{present} & \textbf{present} & avoided & n/a & avoided \\
T9 revenue decoupled from volume & avoided & at-risk & avoided & \textbf{present} & present \\
T10 thin protocol-owned liquidity & n/a & n/a & n/a & at-risk & present \\
T11 consumer equals provider & at-risk & at-risk & avoided & \textbf{present} & present \\
T12 self-signed ranking & at-risk & at-risk & avoided & at-risk & at-risk \\
\bottomrule
\end{tabular}
\end{center}

The sharpest admission is that the self-signed meter, the exact Morpheus defect, is
\code{present} in mord (\code{market/settle.go:101-108}, the doc comment concedes ``a Meter is a
claim, not a proof'' \prov{verified}) and in the on-chain AI mint
(\code{AIMining.submitProof} mints for a proof whose only checks are dedup, a timestamp window,
and a hashcash difficulty test \prov{verified}). mord closes T1 only when the aivm quorum hash
replaces its meter as the settlement input; the fix already exists one directory over
(\code{luxfi/ai/pkg/rewards} ships \code{SlashProvider} and \code{SettleQuorum}), and mord does
not yet import it. We are proposing to sell verification we have not yet made the default in our
own systems. Until at least mord's T1/T2 and \code{lux/ai}'s RPC path are closed, a competent
Morpheus reviewer reading our tree would find we ship the same defect. The punch-list below is
the work that makes the pitch honest.

\subsection{Blocking, before the pitch}

\begin{description}
  \item[B1, DONE 2026-08-29.] \code{lux/ai}'s served RPC was unauthenticated and returned a
    fabricated completion. Fixed: an \code{authMiddleware} (constant-time bearer check, fails
    closed when no key is set) now gates chat, embeddings, and task-submit; the chat handler
    refuses (503/501) instead of fabricating a completion; \code{handleSubmitResult} rejects
    results for unknown tasks, so no unassigned work is credited. Tests in
    \code{cmd/lux-ai/auth\_test.go} pass. The canonical quorum settlement path
    (\code{chains/aivm}) remains authoritative; the placeholder no longer lies about it.
  \item[B2.] vMOR is flat 1:1, not duration-weighted as an earlier note of ours claimed
    (\code{VotesERC20StakedV1.sol:407} mints \code{amount} directly, gating only on one global
    \code{minimumStakingPeriod}). Fix: implement Curve-style duration weighting, or delete the
    claim (already corrected in the token-design note), and set quorum above the owner-directed
    emission share so the vote cannot be captured by emission recipients.
  \item[B3.] The \code{lux/dao} launch Safe is 3-of-5 from one mnemonic; effective control set
    is one (\code{governance.md:22-28}, \code{LUX\_MNEMONIC}). Fix: rotate to independent
    hardware keys held by distinct named people, publish the proofs, and move \code{owner()} of
    the DAO modules to the token-controlled timelock.
  \item[B4.] The \code{lux/standard} production timelock keeps the deployer as sole proposer and
    admin, never renounced (\code{DeployMultiNetwork.s.sol:249-254}). Fix: renounce
    deployer-admin; move the proposer role to the Governor.
  \item[B5.] The reference tenant \code{zoo/vote} points at null contracts
    (\code{config/zoo.json}, governor and factory addresses null). Fix: deploy the governor set
    and fill the config, or stop citing \code{zoo/vote} as a live governance instance. A
    branded front end over no contracts is theater.
\end{description}

\subsection{Required for a real end-state}

\begin{description}
  \item[R1.] The \code{l2-template} chain is centrally sequenced at birth, single-replica MPC,
    KMS, and bridge, with an uncapped-mint admin token (\code{init.sh:240,181,184}). Fix:
    distribute the validator set to independent operators, split MPC/KMS/bridge across distinct
    operators, cap or remove the generated mint, and measure the realized Nakamoto coefficient
    rather than asserting one.
  \item[R2.] Uncapped mint in \code{MOROFT} and the generated \code{Token.sol}
    (\code{MOROFT.sol:50,61}). Fix: for any chain we control, cap or remove the mint. For
    Morpheus, the same fix is ADOPT-03 and SUP-1 on their side.
  \item[R3.] The Lux primary genesis seats about five stakers and five bootstrappers,
    reward-forfeiture only, no principal slashing. Permissionless entry exists in code
    (\code{AddPermissionlessValidatorTx}); the live set is the gap. Fix: open the producer set
    to independent operators and publish the distribution.
  \item[R4.] mord is a single-operator PoC (\code{mord/README.md:5-6}). It is an integrity
    improvement, not a decentralization claim; deploy it under a validator set before calling it
    decentralized. Honest as disclosed.
\end{description}

\subsection{Commitments, the ones that cost us something}

These are not code. They are what we must agree to, or the ``open governance'' claim is theater
regardless of the contracts. We do not hold a controlling share of the vote, and if any reward
flow routes through us we disclose it and cap our weight below a control threshold. We cannot
both direct emission and claim we do not control governance. Permissionless validator entry must
exist in the live set, not only in the code. Keyholders are named people with published proofs.
No retained privileged upgrade path survives once the token-controlled timelock is authoritative.

Every failure above is disclosed in our own code or configs, which by the study's own standard
makes it immaturity, not theater. The discipline is simple: do not make a claim the code does
not yet enforce. Where we are early, we say early. This punch-list is the difference between
shipping that discipline and shipping a better-looking version of what we are asking Morpheus to
leave behind.

%======================================================================
\section{The path}
\label{sec:path}
%======================================================================

Three phases, with a gate between each. A gate is a measured condition, not a date. The
migration is phased and reversible, and the seam is one layer: replace the settlement engine,
keep everything above and below it.

\subsection{Phase 0: on Base, no migration, reversible}

Deploy the four one-transaction fixes (ADOPT-09): route ownership through a timelock (CTL-1),
decode all seven receipt fields and bind the approval address (SET-2, SET-5), and add the
self-deal check (SET-4). Deploy verified settlement as an adjunct to the existing diamond
(ADOPT-02): attestation-gated registration and a quorum check that must confirm before
\code{\_getProviderRewards} pays. Stand up a demand source and attested supply (ADOPT-07): Hanzo
as a first attestable provider through the same config-only backend path any provider uses, with
Enso routing behind \code{api.mor.org}. Providers keep their nodes; consumers see no change.
This phase closes CTL-1, SET-1 through SET-5, ECO-3, and half of SET-3.

\textbf{Gate to Phase 1:} the timelock owns the proxies and the diamond; verified settlement
confirms before payment on a non-trivial share of live sessions; a named demand source is
producing measurable settled volume.

\subsection{Phase 1: governance, finance, bridge}

Deploy governance (ADOPT-03): Safe plus Governor plus timelock with a nonzero delay floor, and a
vMOR vote surface, closing CTL-2 and moving the seat of control off the bare multisig. Route the
buyback's fuel to a fee on verified volume (ADOPT-08). List MOR on the LX order book to deepen a
book that is 88.8\% protocol-owned (ADOPT-06). Deploy LMOR self-repaying credit against the idle
stETH and MOR (ADOPT-04), and the bank rail for fiat-native consumers (ADOPT-05). Bridge MOR
natively across the capital and compute domains and anchor state roots rather than replicating
the schedule by hand (ADOPT-13). MOR the token, its emission, and its capital stay on Ethereum
and Arbitrum; the only change to the token is the buyback's additional source.

\textbf{Gate to Phase 2:} governance parameters sit behind the timelock with a floor above zero;
the fee-fed buyback is live and its dollar flow is published against emission; settled verified
volume is large enough that a few sessions per second is a binding constraint.

\subsection{Phase 2: a chain of its own, thin MOR, emission cut}

Only if the Phase 1 gate is met. Move settlement onto a Lux L1 or L2 (ADOPT-01) that inherits
Quasar PQ finality rather than building consensus, with MOR bridged as gas and mord as the
reference VM; keep MOR thin on top of it (ADOPT-12); lift the throughput ceiling (ADOPT-11).
Run the ICO on \code{lux.exchange} conditional on the regulatory approval \prov{pending}, to
fund the transition and deepen liquidity, not to substitute for demand. Introduce xMOR as a
share-token and cut emission, both of which move the same net-supply inequality from
Section~\ref{sec:mor}. One provenance note we carry from our own audit: ML-KEM is live on Lux
mainnet, while ML-DSA and SLH-DSA are present but not activated, so post-quantum signatures are a
Phase 2 claim, not a current fact.

%======================================================================
\section{What each party gets, and what we must commit to}
\label{sec:parties}
%======================================================================

The migration is a Morpheus governance decision. We cannot execute it, and we should not write
or speak as though the choice is ours. What follows is what each party gets if the community
chooses it, and, stated with equal weight, what we must commit to for it to be real rather than
theater.

\begin{stake}{Morpheus gets}
  \item A token that earns from volume rather than deposit yield, with the emission schedule,
    the holders, and the ticker unchanged.
  \item Settlement that reads the work its receipt already signs, so payment matches the
    whitepaper.
  \item Governance with a delay and a vote where today there is a bare multisig.
  \item Post-quantum consensus, verified inference, threshold custody, and a native bridge,
    inherited rather than built by a five-person team.
\end{stake}

\begin{stake}{Hanzo gets}
  \item The demand and routing layer (\code{api.mor.org} plus Enso), and, running attested
    inference, the largest verified provider on the network.
\end{stake}

\begin{stake}{Lux gets}
  \item The settlement, verification, and transport layer, and the settlement fee for securing
    it.
\end{stake}

\begin{stake}{Zoo gets}
  \item A live governance tenant: MOR holders voting on a real economy through the DAO stack.
\end{stake}

\begin{stake}{What we must commit to, or none of the above is real}
  \item Close our own punch-list first. mord and the on-chain AI mint currently reproduce the
    exact Morpheus settlement defect in our own code; until at least mord's verified path and
    \code{lux/ai}'s RPC path are the default, we are selling verification we have not made our
    own default. B1 is done; the rest are not.
  \item Name a demand source and a volume before claiming the buyback math works. The earning
    model converts demand into token value; it does not manufacture demand. If $V$ is zero, no
    take-rate saves it.
  \item Do not hold a controlling share of the vote, and disclose and cap any reward flow that
    routes through us.
  \item Say early where we are early. Every ``verified live'' in this document means live on
    Lux, Hanzo, or Zoo, not on Morpheus, and every research paper is an argument for adoption,
    not evidence of it.
\end{stake}

The honest obstacle is convincing existing providers to bond stake and accept slashing where
today they bear none, and convincing the multisig to hand parameters to a timelock. Phase 0's
reversibility, and the fact that it turns on real revenue, are the argument that has to carry
the vote. We supply the engine and the proof. The Morpheus community chooses whether to run it.

\end{document}
